\section{Antecedentes}
\label{sec:antecedentes}

La estimación del \emph{Size--of--Wallet} (SioW) y de la \emph{Share--of--Wallet} (SoW) constituye un eje central en la literatura de marketing relacional y analítica de clientes desde finales de la década de 1990. En este cuerpo de trabajo se reconoce que la evaluación del valor económico de un cliente no puede basarse únicamente en el gasto observado en una empresa específica, sino que debe considerar el volumen total de consumo del cliente dentro de una categoría y la proporción de dicho gasto que es capturada por una firma particular. En este marco, el SioW se define como el gasto total del cliente en la categoría, independientemente del proveedor, mientras que la SoW representa la fracción de ese gasto asignada a la empresa focal \cite{du2007size,glady2009predicting}.

Esta distinción resulta fundamental para la gestión estratégica de clientes, ya que un cliente con alto nivel de compras en la empresa focal no necesariamente presenta una alta SoW, y, de manera análoga, un cliente con una SoW elevada puede corresponder a un consumidor de bajo volumen total en la categoría. Trabajos como los de \cite{du2007size} formalizan esta separación conceptual y muestran que la identificación de clientes con alto potencial de crecimiento requiere distinguir entre tamaño absoluto de la billetera y grado de fidelidad relativa hacia la empresa.

En paralelo, diversos estudios han cuestionado la suficiencia de indicadores tradicionales como la satisfacción o el \emph{Net Promoter Score} para explicar el comportamiento económico real de los clientes. En particular, \cite{Keiningham2011} documentan que variaciones modestas en la SoW pueden generar impactos financieros sustancialmente mayores que cambios comparables en medidas de satisfacción declarada, lo que ha impulsado un creciente interés por metodologías cuantitativas orientadas a estimar la SoW de forma directa o indirecta.

Una primera familia de enfoques para estimar SioW y SoW se basa en la recolección directa de información mediante encuestas, paneles de consumidores o diarios de compra. Estos métodos permiten aproximar el gasto total del cliente y su distribución entre marcas o proveedores, proporcionando una referencia cercana al concepto teórico de wallet \cite{du2007size}. No obstante, la literatura reconoce importantes limitaciones prácticas asociadas a estos enfoques, entre las que se incluyen altos costos de implementación, sesgos de memoria y deseabilidad social, baja escalabilidad y dificultades para su integración con bases transaccionales empresariales. Estas limitaciones se acentúan en contextos B2B, donde los procesos de compra suelen ser complejos, agregados y altamente heterogéneos.

Como respuesta a estas restricciones, una segunda línea de investigación ha desarrollado metodologías que infieren el SioW y la SoW utilizando exclusivamente información transaccional y de perfilamiento disponible internamente en las organizaciones. En este grupo se incluyen modelos econométricos y estadísticos que emplean historiales de compra, frecuencias, montos, recencia y características estructurales del cliente para aproximar el comportamiento económico total en la categoría \cite{du2007size}. Estos enfoques reconocen explícitamente que las empresas observan únicamente una fracción del comportamiento real del cliente y buscan inferir la parte no observada a partir de patrones sistemáticos en los datos disponibles.

Dentro de esta familia, una contribución particularmente influyente es la propuesta de \cite{glady2009predicting}, quienes introducen un modelo pogit que separa conceptualmente dos mecanismos latentes: uno asociado al tamaño total de la billetera del cliente en la categoría y otro relacionado con la asignación de dicha billetera hacia la empresa focal. Este enfoque permite modelar de manera explícita la diferencia entre clientes con alta demanda total y clientes con alta propensión a asignar su gasto a una firma específica, utilizando conjuntos diferenciados de covariables transaccionales y de perfilamiento.

La relevancia del modelo de \cite{glady2009predicting} radica en su adecuación a contextos de información parcial, en los cuales no se dispone de la estimación de SoW como input para el modelo. El modelo se construye exclusivamente a partir de transacciones observadas por la empresa focal y permite inferir simultáneamente la SoW y el SioW como variables latentes, lo cual resulta especialmente apropiado para aplicaciones empresariales reales. Además, la separación explícita entre variables transaccionales y variables de perfilamiento proporciona una interpretación económica clara y coherente de los factores que influyen en cada componente del comportamiento del cliente.

De manera complementaria, la literatura en marketing ha desarrollado marcos conceptuales alternativos para explicar la SoW, como la \emph{Wallet Allocation Rule} (WAR) propuesta por \cite{keiningham2015wallet}. Este enfoque enfatiza el carácter multimarca del consumo y plantea que la distribución del gasto entre proveedores depende del posicionamiento relativo de cada marca en la mente del consumidor y del número total de marcas utilizadas. Si bien la WAR se apoya típicamente en información declarativa que no está disponible en la aplicación empírica de este trabajo, su aporte conceptual resulta relevante al reforzar la idea de que la SoW emerge de un proceso de asignación del gasto entre múltiples alternativas, proceso que puede manifestarse indirectamente en patrones transaccionales observables.

Finalmente, la literatura reporta el uso de modelos clásicos de conteo y técnicas de \emph{machine learning} —como árboles de decisión, \emph{random forests} y métodos de \emph{boosting}— para predecir montos de compra o clasificar clientes según niveles de SoW. No obstante, estos enfoques suelen tratar la SoW como una variable observada o exógena y no modelan explícitamente la estructura latente del wallet total del cliente, lo que limita su interpretación en entornos de información incompleta \cite{glady2009predicting}.

En síntesis, el presente trabajo se inscribe en la línea metodológica que busca inferir simultáneamente el tamaño y la asignación de la billetera del cliente a partir de datos transaccionales parciales. Siguiendo el enfoque de \cite{glady2009predicting}, y complementándolo con técnicas de reducción de dimensionalidad y segmentación exploratoria, este estudio adapta dichos modelos al contexto específico de clientes empresariales de una empresa distribuidora de biblias y literatura cristiana, proporcionando una base metodológica sólida para el análisis económico y la toma de decisiones estratégicas en ausencia de información proveniente de encuestas o transacciones con competidores.
